\section{GPU discovery campaign}\label{s:campaign}

The campaign has one logical role: it generates candidate triples $(p,r,k)$.
No lower bound follows from a scanner verdict alone.  A retained candidate
becomes a numerical claim only after the criterion is re-established through
the verification protocol below.  The authenticated coverage-evidence package
supports the coverage and density statements, but is not needed for the
pointwise validity of an accepted certificate.

\subsection{Fused character evaluation}\label{ss:b7}

The scan kernel exploits the fact that the additive order of positions is free:
we iterate $x=1,2,\dots,p-1$ directly and evaluate the colour of $x$ on the fly
by modular exponentiation, fused with the run-length bookkeeping. Nothing is
sorted and no per-element value is materialized in memory. A single pass
evaluates $y=x^{(p-1)/m}$ with $m=\operatorname{lcm}\{\,r\le 9 : r \mid p-1\,\}
\le 2520$; a table lookup of $y$ among the $m$-th roots of unity gives
$\operatorname{ind}(x) \bmod m$, hence the colour for every modulus $r\le 9$
simultaneously. Runs are summarized per segment by (prefix, suffix, best
interior) triples combined with an associative, order-respecting reduction.

A second symmetry halves the work. Since $c(p-x)=c(x)+c(-1)$, the second half
of the sequence is the reversal of the first half shifted by the constant
$c(-1)$; it suffices to scan $x=1,\dots,(p-1)/2$ and to splice the two halves
at the center. We refer to this half-interval variant as the mirror scan.

\subsection{Redundant and separate-source validation paths}\label{ss:validation}

The retained CUDA qualification compares the mirror, full-interval,
sort-based and CPU-walk paths on small primes, tests the boundary condition
by two constructions, and exercises the high-range modular arithmetic.
Its raw logs and source identities are preserved with the qualification
manifest. The current CPU suite also tests the order-respecting run reduction.
The historical development exposed why this matters: run summaries compose
associatively but not commutatively, so a reduction that reorders them can
produce an incorrect profile. The archived scanner uses the repaired
ordered reduction. The table below identifies the source and arithmetic
shared by each path; Section~\ref{s:data} states the replication status.
\begin{table}[ht]
\centering\scriptsize
\begin{tabular}{P{2.3cm} P{1.7cm} P{2.4cm} P{2.4cm} P{3.0cm}}
\toprule
path & language & modular arithmetic & class/run construction & shared components \\
\midrule
mirror V2a & CUDA C++ & GPU Montgomery & character, half interval & same source and helpers as full B7 \\
full B7 & CUDA C++ & GPU Montgomery & character, full interval & same source and helpers as V2a \\
sort-based GPU & CUDA C++ & GPU Montgomery & sort then run scan & earlier GPU design; shares GPU Montgomery arithmetic \\
CPU walk & C++ host & native 64-bit walk & discrete logarithm, sequential & same executable and host primitive-root code \\
CPU criterion & C++ & separate 128-bit path & separate walk and (B$^*$) & separate source and binary \\
\texttt{verify\_claim} & C++ & 128-bit exponentiation & roots of unity and full criterion & strongest source/arithmetic separation \\
high-range witness & C & no Montgomery & full sequential walk & separate arithmetic; tests (a), not (B$^*$) \\
\bottomrule
\end{tabular}
\caption{Dependency matrix for the verification paths.  V2a and B7 are
correlated redundant paths.  The principal source- or arithmetic-separated
checks are the separate CPU criterion, \texttt{verify\_claim}, and, for the
high-range arithmetic defect, the Montgomery-free witness.}
\label{tab:dependencies}
\end{table}

\subsection{Claim-verification protocol}\label{ss:protocol}

The certificate registry contains exactly thirteen direct Rabung
claims: nine three-colour triples for $17\le k\le25$ and four two-colour
triples for $25\le k\le28$.  For each triple $(p,r,k)$ the release protocol
requires:
\begin{enumerate}
\item two executions of the structured triplet \{mirror scan, full scan,
  CPU walk\}, with internal profile equality and equality between executions;
\item two acceptances by the separate full CPU criterion;
\item acceptance by the stand-alone \texttt{verify\_claim}; and
\item for the two primes above the historical Montgomery-safety threshold, an
  additional full-interval witness that uses no Montgomery arithmetic.
\end{enumerate}
The parser is fail-closed: it checks the process status, the exact $p,r,k$,
all profile equalities, (a), (B$^*$), and the final verdict, and rejects any
output containing \texttt{DIFF}, \texttt{FAIL}, \texttt{ECHEC}, or a negative
verdict.  Its digest binds the claim identity, profiles, criterion values, and
final verdict.  The manifest separately records source and binary hashes,
compiler version, exact commands, and hashes of every raw output.  The final
Markdown and \LaTeX{} claim tables are generated only from the exact set of
thirteen accepted manifests.

This rule is stricter than the historical transcript, whose max-run checksum
did not bind $p$ and whose pasted \texttt{verify\_claim} outputs were sometimes
truncated. Those files remain historical records. The targeted PC run
supplied the thirteen complete manifests, recording $67$ required executions: five per claim plus the
two additional Montgomery-free high-range witnesses.

\paragraph{Current CUDA qualification.}
The recorded qualification builds the committed scanner for \texttt{sm\_120}
and passes the complete qualification suite on the target RTX~5090.  Its
manifest binds the source and binary hashes, compiler, driver, GPU identity,
commands, logs, exit status, and clean pre-build state.  This qualifies the
discovery implementation and hardware path; it does not replace any Rabung
certificate check.

\subsection{Campaign execution}\label{ss:execution}

The production interval was $[970\,000\,000,2\,000\,000\,000)$, partitioned
into $515$ contiguous chunks of width $2\times10^6$. Each prime was evaluated
for the active colour counts $r\le9$ dividing $p-1$. The checkpoint accounts
for every chunk; $514$ also have matching journal records, while chunk~7 is
recovered from the checkpoint alone. The surviving outputs contain $498$
per-chunk histograms and one aggregate re-scan of the first $17$ chunks.

A separately written CPU sieve and the scanner produce byte-identical
ordered prime streams on all $515$ chunks, containing $48\,823\,489$ primes.
This checks the enumeration underlying the coverage statement. It does not
replace the pointwise criterion tests for a certificate. Four historical
scanner fingerprints were used, with no change to the $(3,17)$ target
predicate; the complete holdout segment uses the final fingerprint.
Appendix~\ref{s:campaign-audit} gives the revision table, the distinction
between histogram and ordered-stream checks, and the recovery of the early
aggregate evidence.

The GPU method made this finite survey practicable, but the retained material
does not constitute a controlled throughput study with a common repetition
and telemetry protocol. We therefore report the algorithm and coverage,
without using historical kernel rates as publication benchmarks. None of the
numerical lower bounds depends on a performance claim.
